% !TeX root = ../algebra-1.tex
\chapter{Orthogonality}\label{chap:algebra-orthogonality}

\section{Conjugation, inner products and norms}

Complex conjugation is applied entry by entry to vectors and matrices:
\[
 \overline{A+B}=\overline A+\overline B,\qquad
 \overline{AB}=\overline A\,\overline B,\qquad
 \overline{A^{\mathsf T}}=(\overline A)^{\mathsf T}.
\]
The \emph{conjugate transpose}, or adjoint, is $A^*=\overline A^{\mathsf T}$.
Transposition reverses products, so $(AB)^*=B^*A^*$.
For real matrices, $A^*=A^{\mathsf T}$.

\begin{definition}[The source's inner product convention]
For $u,v\in F^n$, define
\[
 \boxed{u\cdot v=u^{\mathsf T}\overline v
       =\sum_{j=1}^n u_j\overline{v_j}.}
\]
The \emph{norm}, or length, is
\[
 \|u\|=\sqrt{u\cdot u}=\sqrt{|u_1|^2+\cdots+|u_n|^2}.
\]
\end{definition}

For $z=x+iy$, $|z|^2=z\overline z=x^2+y^2$, whereas
$|z|=\sqrt{x^2+y^2}$.
Thus $u\cdot u$ is nonnegative real, and is zero exactly when $u=0_n$.
For real vectors the inner product is the usual dot product.

\begin{proposition}[Inner product rules]
For compatible vectors and matrices,
\begin{align*}
 u\cdot v&=\overline{v\cdot u},\\
 (au+bw)\cdot v&=a(u\cdot v)+b(w\cdot v),\\
 u\cdot(av+bw)&=\overline a(u\cdot v)+\overline b(u\cdot w),\\
 (Au)\cdot v&=u\cdot(A^*v).
\end{align*}
\end{proposition}
\begin{proof}
The first three follow by expanding the coordinate sums. For the last,
\[
 (Au)\cdot v=u^{\mathsf T}A^{\mathsf T}\overline v
 =u^{\mathsf T}\overline{A^*v}=u\cdot(A^*v).
\]
\end{proof}

\begin{remark}
The inner product is linear in the \textbf{first} argument and
conjugate-linear in the second: a scalar leaving the second argument
must be conjugated. Other books may use the opposite convention.
All projection and Gram--Schmidt coefficients below follow our convention.
\end{remark}

\begin{example}[Checking the complex arithmetic]
Let $u=(2+i,-1-i)^{\mathsf T}$, $v=(1,1)^{\mathsf T}$ and
$w=(-1,2)^{\mathsf T}$. Then
\begin{align*}
 u\cdot u&=|2+i|^2+|-1-i|^2=5+2=7,\\
 u\cdot v&=(2+i)+(-1-i)=1,\\
 v\cdot w&=-1+2=1.
\end{align*}
Thus $\|u\|=\sqrt7$. The value of $u\cdot v$ corrects the source's arithmetic.
\end{example}

\section{Orthogonality, Cauchy--Schwarz and angles}

\begin{definition}
Vectors $u,v$ are \emph{orthogonal} if $u\cdot v=0$.
A sequence is orthogonal if each pair of distinct vectors is orthogonal.
\end{definition}

Conjugate symmetry makes the relation symmetric.
For orthogonal $u,v$, expansion gives Pythagoras's identity:
\[
 \|u+v\|^2=(u+v)\cdot(u+v)=\|u\|^2+\|v\|^2.
\]

\begin{theorem}[Cauchy--Schwarz inequality]
For $u,v\in F^n$,
\[
 \boxed{|u\cdot v|\leq\|u\|\,\|v\|.}
\]
Equality holds if and only if $u,v$ are linearly dependent.
\end{theorem}
\begin{proof}
If $v=0_n$, both sides are zero and the pair is dependent.
Otherwise put $c=(u\cdot v)/(v\cdot v)$ and $w=u-cv$.
Then $w\cdot v=0$, so Pythagoras applied to $u=cv+w$ gives
\[
 \|u\|^2=|c|^2\|v\|^2+\|w\|^2
 \geq\frac{|u\cdot v|^2}{\|v\|^2}.
\]
Multiply by $\|v\|^2$ and take square roots.
Equality holds exactly when $w=0$, that is, $u=cv$.
\end{proof}

\begin{definition}[Angles between real vectors]
For nonzero $u,v\in\mathbb R^n$, their angle $\theta\in[0,\pi]$ satisfies
\[
 \cos\theta=\frac{u\cdot v}{\|u\|\,\|v\|}.
\]
\end{definition}

Cauchy--Schwarz ensures the quotient is in $[-1,1]$.
Real nonzero vectors are orthogonal exactly when their angle is $\pi/2$.
For complex vectors the quotient may be nonreal, so this ordinary
real-angle definition cannot simply be used unchanged.
Orthogonality itself remains well-defined over $\mathbb C$.

\begin{proposition}
An orthogonal sequence of nonzero vectors is linearly independent.
\end{proposition}
\begin{proof}
If $\sum_j a_jv_j=0$, take the inner product with $v_k$ on the right.
This gives $a_k\|v_k\|^2=0$. Since $\|v_k\|^2>0$, every $a_k=0$.
\end{proof}

\section{Orthonormal bases}

\begin{definition}
A sequence $q_1,\ldots,q_k$ is \emph{orthonormal} if
\[
 q_i\cdot q_j=\begin{cases}1&i=j,\\0&i\ne j.\end{cases}
\]
Thus its vectors are mutually orthogonal and have norm one.
\end{definition}

An orthogonal nonzero sequence becomes orthonormal by replacing each
$v_j$ with $v_j/\|v_j\|$, which does not change its span.

\begin{theorem}[Coordinates in an orthonormal basis]
If $\mathcal B=(q_1,\ldots,q_d)$ is an orthonormal basis of $U$, then
for $v\in U$,
\[
 \boxed{v=\sum_{j=1}^d(v\cdot q_j)q_j,\qquad
 [v]_{\mathcal B}=(v\cdot q_1,\ldots,v\cdot q_d)^{\mathsf T}.}
\]
\end{theorem}
\begin{proof}
Write $v=\sum_j c_jq_j$ and take the inner product with $q_k$ on the
right. Linearity and orthonormality give $v\cdot q_k=c_k$.
\end{proof}

\begin{example}
An orthonormal basis of $\mathbb R^2$ is
\[
 q_1=\frac1{\sqrt2}(1,1)^{\mathsf T},\qquad
 q_2=\frac1{\sqrt2}(1,-1)^{\mathsf T}.
\]
For $v=(1,3)^{\mathsf T}$, the coordinates are
$v\cdot q_1=2\sqrt2$ and $v\cdot q_2=-\sqrt2$.
Hence $v=2\sqrt2\,q_1-\sqrt2\,q_2$.
\end{example}

\section{Orthogonal projection}

\begin{definition}
The projection of $v$ onto the direction of a nonzero vector $u$ is
\[
 \boxed{\operatorname{proj}_u(v)=\frac{v\cdot u}{u\cdot u}\,u.}
\]
\end{definition}

The projection belongs to $\spanop(u)$, and its residual is orthogonal to $u$:
\[
 \left(v-\frac{v\cdot u}{u\cdot u}u\right)\cdot u=0.
\]
It is the unique vector in $\spanop(u)$ with this property:
if $v-cu$ is orthogonal to $u$, the equation forces
$c=(v\cdot u)/(u\cdot u)$.

\begin{center}
% Drawing only: edit the chapter text in chapters/algebra-05-orthogonality.tex.
\begin{tikzpicture}[>=Stealth,scale=1.05]
 \coordinate (O) at (0,0);
 \coordinate (P) at (3,0);
 \coordinate (V) at (3,1.8);
 \draw[->,muted] (-0.4,0)--(4.4,0) node[right] {$\spanop(u)$};
 \draw[->,thick,ink] (O)--(V) node[midway,above left] {$v$};
 \draw[->,very thick,accent] (O)--(P)
   node[midway,below] {$p=\operatorname{proj}_u(v)$};
 \draw[->,dashed,ink] (P)--(V) node[midway,right] {$v-p$};
 \draw (2.8,0)--(2.8,0.2)--(3,0.2);
 \node[below left] at (O) {$0$};
\end{tikzpicture}

\end{center}

For real vectors with angle $\theta$, the signed component along
$u/\|u\|$ is $\|v\|\cos\theta$.
The projection's length is $\|v\|\,|\cos\theta|$;
the absolute value matters for an obtuse angle.

\subsection{Projection onto a subspace}

For an orthonormal basis $q_1,\ldots,q_d$ of $U$, define
\[
 p=\operatorname{proj}_U(v)=\sum_{j=1}^d(v\cdot q_j)q_j.
\]
Then $p\in U$ and $v-p$ is orthogonal to every vector of $U$.
These properties determine $p$ uniquely, independently of the basis:
if $p,p'$ both work, their difference belongs to $U$ and is orthogonal
to $U$, hence to itself. It has norm zero, so $p=p'$.

\begin{proposition}[Closest point interpretation; additional clarification]
The projection $p$ is the unique vector in $U$ closest to $v$.
\end{proposition}
\begin{proof}
For $y\in U$, $v-y=(v-p)+(p-y)$ is an orthogonal decomposition.
Thus $\|v-y\|^2=\|v-p\|^2+\|p-y\|^2\geq\|v-p\|^2$,
with equality exactly when $y=p$.
\end{proof}

\section{The Gram--Schmidt process}

Given independent $b_1,\ldots,b_k$, subtract projections onto the
previously constructed directions:
\begin{align*}
 v_1&=b_1,\\
 v_2&=b_2-\frac{b_2\cdot v_1}{v_1\cdot v_1}v_1,\\
 v_j&=b_j-\sum_{i=1}^{j-1}
             \frac{b_j\cdot v_i}{v_i\cdot v_i}v_i
             \qquad(2\leq j\leq k).
\end{align*}
Finally set $q_j=v_j/\|v_j\|$.

\begin{theorem}[Gram--Schmidt]
Each $v_j$ is nonzero, the $v_j$ are orthogonal, and
\[
 \spanop(v_1,\ldots,v_j)=\spanop(b_1,\ldots,b_j)
 \quad(1\leq j\leq k).
\]
The $q_j$ are orthonormal. Every subspace of $F^n$ therefore has an
orthonormal basis.
\end{theorem}
\begin{proof}
Induct on $j$. Independence makes $v_1=b_1$ nonzero.
Assume the assertions hold through $j-1$.
For $\ell<j$, take the inner product with $v_\ell$ on the right:
\[
 v_j\cdot v_\ell=b_j\cdot v_\ell
 -\frac{b_j\cdot v_\ell}{v_\ell\cdot v_\ell}
 (v_\ell\cdot v_\ell)=0.
\]
All other terms vanish by orthogonality.
The definition subtracts a vector in the preceding span, so adjoining
$v_j$ or adjoining $b_j$ gives the same enlarged span.
If $v_j=0$, then $b_j$ belongs to the span of its predecessors, contrary
to independence. All denominators are therefore positive and
normalisation is valid. Starting with a basis of a subspace proves the
last assertion.
\end{proof}

\begin{remark}
Over $\mathbb C$, the numerator is $b_j\cdot v_i$, not $v_i\cdot b_j$,
under our convention. They are conjugates and need not be equal.
\end{remark}

\begin{example}[A complete three-dimensional calculation]
Start with
\[
 b_1=\begin{pmatrix}1\\2\\3\end{pmatrix},\quad
 b_2=\begin{pmatrix}4\\5\\6\end{pmatrix},\quad
 b_3=\begin{pmatrix}7\\8\\10\end{pmatrix}.
\]
Their determinant is $-3$, so they are independent.
As $\|b_1\|^2=14$ and $b_2\cdot b_1=32$,
\[
 v_1=b_1,\qquad
 v_2=b_2-\frac{16}{7}b_1
     =\frac37\begin{pmatrix}4\\1\\-2\end{pmatrix}.
\]
Projection onto a direction is unchanged by rescaling its generator,
so use $w_2=(4,1,-2)^{\mathsf T}$ at the next step.
We have $b_3\cdot b_1=53$, $b_3\cdot w_2=16$ and $\|w_2\|^2=21$.
Thus
\begin{align*}
 v_3&=b_3-\frac{53}{14}b_1-\frac{16}{21}w_2\\
 &=\frac16\begin{pmatrix}1\\-2\\1\end{pmatrix}.
\end{align*}
After normalisation,
\[
 \boxed{
 q_1=\frac1{\sqrt{14}}\begin{pmatrix}1\\2\\3\end{pmatrix},\quad
 q_2=\frac1{\sqrt{21}}\begin{pmatrix}4\\1\\-2\end{pmatrix},\quad
 q_3=\frac1{\sqrt6}\begin{pmatrix}1\\-2\\1\end{pmatrix}.}
\]
Their pairwise dot products vanish and each has norm one.
\end{example}

\method{Keep the orthogonal vectors unnormalised during subtraction,
then normalise at the end to avoid unnecessary square roots.
A zero residual indicates dependence in the input sequence.}

\section{Orthogonal and unitary matrices}

\begin{definition}
A real square matrix $Q$ is \emph{orthogonal} if $Q^{\mathsf T}Q=I_n$.
A complex square matrix $Q$ is \emph{unitary} if $Q^*Q=I_n$.
\end{definition}

For an orthogonal matrix, $Q^{-1}=Q^{\mathsf T}$; for a unitary matrix,
$Q^{-1}=Q^*$.
A square matrix with a left inverse is invertible, by the
invertibility theorem.

\begin{proposition}
A real square matrix is orthogonal exactly when its columns are an
orthonormal basis of $\mathbb R^n$.
The complex analogue holds with unitary matrices and $\mathbb C^n$.
\end{proposition}
\begin{proof}
For real columns $q_i$, entry $(i,j)$ of $Q^{\mathsf T}Q$ is $q_i\cdot q_j$.
For complex columns, entry $(i,j)$ of $Q^*Q$ is $q_j\cdot q_i$ under
our convention. It equals $\delta_{ij}$ exactly when the columns are
orthonormal. There are $n$ independent columns, so they form a basis.
\end{proof}

Orthogonal and unitary matrices preserve inner products:
\[
 (Qu)\cdot(Qv)=u\cdot(Q^*Qv)=u\cdot v.
\]
Thus they preserve lengths and orthogonality.
Real orthogonal matrices also preserve real angles.
Taking determinants in $Q^{\mathsf T}Q=I_n$ gives
$(\det Q)^2=1$, so $\det Q=1$ or $-1$.

\begin{remark}
An orthogonal \emph{matrix} has orthonormal columns, including length one.
Merely having orthogonal nonzero columns is insufficient:
$2I_n$ has orthogonal columns but is not an orthogonal matrix.
\end{remark}

\section{Real symmetric matrices and the spectral theorem}

\begin{definition}
A real matrix is \emph{symmetric} if $A^{\mathsf T}=A$.
A complex matrix is \emph{Hermitian} if $A^*=A$.
\end{definition}

\begin{lemma}
Every complex eigenvalue of a real symmetric matrix is real.
Eigenvectors for distinct eigenvalues are orthogonal.
\end{lemma}
\begin{proof}
Let $Av=\lambda v$ with $v\in\mathbb C^n$ nonzero. Since $A^*=A$,
\[
 \lambda\|v\|^2=(Av)\cdot v
 =v\cdot(Av)=\overline\lambda\|v\|^2.
\]
The positive norm squared gives $\lambda=\overline\lambda$.
If $Au=\lambda u$ and $Av=\mu v$, then
\[
 \lambda(u\cdot v)=(Au)\cdot v=u\cdot(Av)=\mu(u\cdot v).
\]
Distinct real $\lambda,\mu$ force $u\cdot v=0$.
\end{proof}

\begin{lemma}[A real eigenvector can be chosen]
A real symmetric matrix of positive size has a nonzero real eigenvector.
\end{lemma}
\begin{proof}
Its characteristic polynomial has a complex root $\lambda$; the
previous lemma makes $\lambda$ real.
Write a complex eigenvector as $v=x+iy$ with $x,y\in\mathbb R^n$.
Because $A,\lambda$ are real, $Av=\lambda v$ gives
$Ax=\lambda x$ and $Ay=\lambda y$.
At least one of $x,y$ is nonzero, giving a real eigenvector.
\end{proof}

\begin{theorem}[Real spectral theorem; source Theorem 5.22, strengthened form]
Every real symmetric matrix $A$ has an orthonormal basis of real
eigenvectors. Equivalently, a real orthogonal matrix $Q$ exists with
\[
 \boxed{Q^{\mathsf T}AQ=D,\qquad A=QDQ^{\mathsf T},}
\]
where $D$ is real and diagonal. In particular, $A$ is diagonalisable
over $\mathbb R$.
\end{theorem}
\begin{proof}
Induct on $n$, with $n=1$ immediate.
Choose a real unit eigenvector $q_1$ with eigenvalue $\lambda$.
Extend it to a basis and apply Gram--Schmidt.
Since $q_1$ is already a unit vector, it remains the first vector.
This gives an orthogonal $Q_0=[\,q_1\ q_2\ \cdots\ q_n\,]$.

The first column of $Q_0^{\mathsf T}AQ_0$ is
$(\lambda,0,\ldots,0)^{\mathsf T}$.
Symmetry forces the same zeros in the first row, so
\[
 Q_0^{\mathsf T}AQ_0=\begin{pmatrix}\lambda&0\\0&B\end{pmatrix},
\]
where $B$ is real symmetric of size $n-1$.
By induction, an orthogonal $Q_1$ makes $Q_1^{\mathsf T}BQ_1=D_1$ diagonal.
Then
\[
 Q=Q_0\begin{pmatrix}1&0\\0&Q_1\end{pmatrix}
\]
is orthogonal and $Q^{\mathsf T}AQ=\diag(\lambda,D_1)$.
Its columns are the required eigenvector basis.
\end{proof}

\begin{example}
For $A=\begin{pmatrix}2&1\\1&2\end{pmatrix}$, take
$Q=\frac1{\sqrt2}\begin{pmatrix}1&1\\1&-1\end{pmatrix}$.
Its columns are orthonormal eigenvectors and
$Q^{\mathsf T}AQ=\diag(3,1)$.
This is the diagonalisation from Chapter~\ref{chap:algebra-linear-maps},
now with an orthogonal change of basis.
\end{example}

\begin{remark}[Complex symmetric is different from Hermitian]
Over $\mathbb C$, $A^{\mathsf T}=A$ alone does not imply diagonalisability.
For example,
\[
 A=\begin{pmatrix}1&i\\i&-1\end{pmatrix}
\]
is symmetric and nonzero but $A^2=0$.
Any eigenvalue must be zero; diagonalisation would therefore force
$A=0$, a contradiction.
The complex spectral theorem requires $A^*=A$:
a Hermitian matrix has real eigenvalues and a unitary diagonalisation
$Q^*AQ=D$. This is the Hermitian analogue mentioned in the source.
\end{remark}
